\scanpage{6}
\deleted{\unsure{$G$-$\underline{\mathrm F}$ $G$-$\underline{\mathrm F}_p$}($G$)-Modules} Modules de représentations de $G$ sur $\underline{\mathrm F}_p$. Alors les termes $(\mathrm{Coker}\,u_2)_{ss}$ et $(\mathrm{Ker}\,u_2)_{ss}$ peuvent s'interpréter aussi comme les duals de $(\mathrm{Ker}\,u''_2)_{ss}$ et de $(\mathrm{Coker}\,u''_2)_{ss}$, $u''_2$ défini dans (8.3.2). Enfin \deleted{$(\mathrm{Ker}\,u_1$} $(\mathrm{Coker}\,u_1)_{ss}$ et $(\mathrm{Ker}\,u_1)_{ss}$ ont aussi une interprétation simple, grâce au

\uline{Lemme} 8.4.\deleted{\unsure{2}}\handwritten{3}. Soit $X$ un schéma propre sur un corps $k$, de car. $p\handwritten{>}0$. Alors pour tout $i\handwritten{\in}\underline Z$, l'homomorphisme d'augmentation \deleted{\unsure{$H$}}
\[
H^i_{\rm DR}(X)\handwritten{\longrightarrow}H^i(X,\underline O_X)
\]
induit un isomorphisme sur les parties semi-simples (pour l'opération $p$-linéaire de Frobenius).

Ceci résulte aussitôt de la suite spectrale de Hodge
\[
E^{a,b}_1=H^b(X,\underline\Omega^a_{X/k})\handwritten{\Longrightarrow}H^{\bullet}_{\rm DR}(X)\;,
\]
compte tenu que $\mathrm{Frob}$ opère \deleted{semi} $p$-linéairement sur cette suite spectrale \inserted{et} qu'il est nul sur $E_1^{a,b}$ pour $a\handwritten{\ne}0$. Alors les $(E_1^{0,b})\deleted{=}_{ss}=H^b(X,\underline O_X)_{ss}$ sont des cycles universels pour les $d_r$, et on gagne ...
% margin note written sideways
\margin{\handwritten{Cor. 8.4.\uline{$4\unsure{\mathrm{bis}|\mathrm{ter}}$} (utilisant 8.1.\unsure{?|2}) les résultats de la feuille jointe (marquée 8.4.4).}}

Mettant ensemble les renseignements obtenus, \deleted{et en même temps, pour des raisons de commodité, transformant par dualité de Poincaré l'homomorphisme $u''_2$ de (8.3.2), de sorte qu'on trouve (oubli dans 8.3 !)}
\[
\deleted{u_2\quad\text{est can. isomorphe à}}
\tag*{\deleted{(8.4.2)}}
\]
on trouve:

``\uline{Théorème}'' 8.5. En admettant les\handwritten{\,}ingrédients précédents pour la cohomologie cristalline (\inserted{formule des} coefficients universels, dualité de Poincaré, théor\deleted{\unsure{ème}}\handwritten{ème} de Lefschetz \handwritten{``faible''}\deleted{, y compris pour les groupes de torsion} \inserted{\handwritten{sans négliger la torsion}}) on obtient une suite exacte\handwritten{,} relative à une variété lisse projective connexe $X\handwritten{\subset}\mathrm P^r_k$ \inserted{de di\deleted{\unsure{p}}mension $n$} sur un corps parfait, de car. $p\handwritten{>}0$, et une section hyperplane (ou par une hypersurface \inserted{\handwritten{de degré $d\geq1$}}, cela s'y ramène ...) $Y$ lisse de dimension $n-1$:
\[
\begin{aligned}
&\deleted{(8.5.1)\quad 0\quad (\mathrm{Ker}\,\unsure{v})\quad(\unsure{\mathrm{Coker}}\,\unsure{w})\unsure{'}}\\
&\deleted{(E(Y)\otimes_W k)_{ss}\quad\deleted{\mathrm{Coker}\,v}\unsure{E_{\rm DR}}}
\end{aligned}
\]
% archive sheet number: 5

